\documentclass[11pt]{letter}

\usepackage{geometry}
\geometry{margin=1in}
\usepackage{setspace}
\usepackage{hyperref}

\signature{Decory Edwards}
\address{Decory Edwards \\ Johns Hopkins University \\ Department of Economics}

\begin{document}

\begin{letter}{Hiring Committee \\ Digital Capital Advisors \\ New York, NY}

\opening{Dear Hiring Committee,}

I am writing to apply for the 2026 Full-Time Technology Investment Banking Analyst position at Digital Capital Advisors. I am currently a PhD candidate in Economics at Johns Hopkins University. My training combines mathematical reasoning, financial and computational modeling, and data driven empirical analysis. I am motivated by problems that require precise logic, careful implementation, and constant validation against data. DCA’s focus on technology driven transactions and its entrepreneurial deal environment align closely with how I approach analytical work.

My research agenda is at the intersection of wealth inequality and macroeconomics. In my job market paper, I build and estimate a structural heterogeneous agent model with returns that differ across households. The core contribution is to show how heterogeneity in returns and income risk jointly shape the wealth distribution and consumption behavior. Relative to closely related work, my framework performs better at matching the lower and middle portions of the wealth distribution. This difference matters quantitatively. Models that focus primarily on returns heterogeneity can fit the top of the wealth distribution closely but tend to generate too much wealth among the bottom half of households. In my framework, income uncertainty generates a precautionary saving motive that produces genuinely low wealth households. This leads to a realistic distribution of marginal propensities to consume and aggregate consumption responses that align with empirical estimates.

A key feature of my research is the heavy use of modeling and computation to analyze problems that cannot be handled analytically. I build structured models, evaluate outcomes across scenarios, and assess sensitivity to assumptions. I work with large datasets, construct new variables from raw information, and present results clearly. This work requires careful numerical reasoning, attention to detail, and disciplined validation, skills that are directly applicable to valuation work, transaction analysis, and deal execution.

I am particularly interested in Digital Capital Advisors because of its focus on technology and digital media companies and its hands-on deal execution model. The opportunity to work closely with senior bankers on M\&A transactions and capital raises in a small team environment is especially appealing. I value settings where analysts take ownership of their work and where analytical rigor directly informs decisions.

My economics training has provided a strong foundation in probability, optimization, and financial reasoning. I am comfortable translating theory into practical models and working through complex quantitative problems under tight deadlines. I have extensive experience using Excel and quantitative tools to build transparent and reliable models. I also regularly produce clear written materials and presentations to communicate results. These skills are integral to my research process and would be central to my contributions as an analyst at DCA.

I am also enthusiastic about the prospect of living and working in New York City. I value in person collaboration and the intensity of a fast paced investment banking environment. I am comfortable committing fully to an on-site role and see this position as an opportunity to apply my analytical training in a demanding and execution focused setting.

I would welcome the opportunity to contribute my skills and perspective to Digital Capital Advisors. Thank you for your time and consideration.

\closing{Sincerely,}

\end{letter}
\end{document}
